\documentclass[12pt,a4paper]{article}

\usepackage{amsmath,amssymb,amsthm}
\usepackage{geometry}
\usepackage{hyperref}
\usepackage{booktabs}
\usepackage{graphicx}
\usepackage{xcolor}
\usepackage{url}
\usepackage{cite}
\usepackage{caption}
\usepackage{array}
\usepackage{enumitem}
\usepackage{tabularx}
\geometry{margin=2.5cm}

\newtheorem{theorem}{Theorem}[section]
\newtheorem{lemma}[theorem]{Lemma}
\newtheorem{corollary}[theorem]{Corollary}
\newtheorem{proposition}[theorem]{Proposition}
\newtheorem{conjecture}[theorem]{Conjecture}
\newtheorem{definition}{Definition}[section]
\newtheorem{assumption}{Assumption}[section]
\newtheorem{remark}{Remark}[section]

\definecolor{navy}{HTML}{1a2744}
\definecolor{gold}{HTML}{c9a84c}
\definecolor{teal}{HTML}{1a6b5a}

\hypersetup{
  colorlinks=true,
  linkcolor=navy,
  citecolor=teal,
  urlcolor=teal,
}

\graphicspath{{./figures/}{./}}

% ------------------------------------------------------------------
\title{%
  \textbf{Principia Orthogona, Volume~I}\\[0.3em]
  \textbf{The Mathematics of Generative Transitions}\\[0.5em]
  \large Second Edition --- April 2026%
}

\author{%
  Pablo Nogueira Grossi\\
  G6 LLC, Newark, New Jersey, USA\\
  ORCID: 0009-0000-6496-2186\\
  \texttt{pgrossi888@outlook.com} \quad \texttt{g6llc@proton.me}\\[0.3em]
  \small AXLE repository: \url{https://github.com/TOTOGT/AXLE}\\
  \small DM3-lab: \url{https://github.com/TOTOGT/DM3-lab}\\
  \small Zenodo: \href{https://doi.org/10.5281/zenodo.19117400}{10.5281/zenodo.19117400}\\
  \small ISBN 979-8-9954416-0-1%
}

\date{April 2026}
% ------------------------------------------------------------------

\begin{document}
\maketitle

% ------------------------------------------------------------------
\begin{abstract}
This volume develops a unified mathematical framework for generative
transitions: localised geometric events in which a trajectory undergoes
compression, curvature intensification, loss of injectivity, and
stabilisation, governed by the operator sequence
$G = U \circ F \circ K \circ C$.
The framework rests on six minimal assumptions (manifold structure,
trajectory regularity, compression feasibility, curvature threshold,
folding well-posedness, Morse stability functional) and produces:
constructive operator definitions with explicit formulas; five
structural theorems including existence, non-commutativity, and finite
branching; seven analytical invariants; four normal forms; a singularity
classification restricted to the Whitney $A_1$--$A_3$ hierarchy;
a free-discontinuity variational principle; and a symplectic
Hamiltonian structure with a distributional generator at the fold.
The second edition extends the framework by identifying a fifth
operator $E$ (Generative Time Circuit) whose action variable $z$
satisfies $\dot{z} \geq 0$ monotonically along trajectories near the
limit cycle, accumulating the dissipative cost of each generative cycle.
A term-by-term structural correspondence is established between the
five-operator chain $C \to K \to F \to U \to E$ and Perelman's proof
of the Poincar\'e conjecture via Ricci flow with surgery
\cite{Perelman2002,Perelman2003a,Perelman2003b}, stated as
Conjecture~\ref{conj:perelman}.
The dimensional threshold $N = 3$ is identified as the minimum
dimension for non-trivial contact geometry, connecting it to the
threshold $c = 3$ in the Collatz map
(Conjecture~\ref{conj:threshold}).
Theorems~A--D (existence, closure, contact normal form, structural
stability) are machine-checked in Lean~4 via the AXLE repository
\cite{AXLE2026} with zero axioms beyond Mathlib4.
\end{abstract}

\textbf{Keywords:} generative transitions; contact geometry;
operator sequence; Whitney fold; singularity theory; variational
mechanics; symplectic geometry; Ricci flow; Perelman conjecture;
Lean~4 formal verification; dimensional threshold

\textbf{MSC codes:} 37C25, 37G10, 53D10, 57M27, 58K05, 70H05, 47H10

% ------------------------------------------------------------------
\tableofcontents
\bigskip
% ------------------------------------------------------------------

% ==================================================================
\section{Introduction}
\label{sec:intro}
% ==================================================================

A generative transition is what happens when a system crosses a
threshold it cannot uncross.
The cell that begins to digest itself under nutrient stress.
The star that ignites helium when its hydrogen is exhausted.
The elastic rod that buckles under axial load.
The 3-manifold that develops a geometric singularity under Ricci flow.
In each case, a compression drives the system toward a critical point,
a curvature instability makes the approach inevitable, a fold commits
the system irreversibly to a new branch, and an unfolding establishes
the new stable state.

This volume formalises that sequence as four operators acting on
trajectories in a Riemannian manifold:
\[
  C \;\longrightarrow\; K \;\longrightarrow\; F
  \;\longrightarrow\; U.
\]
The second edition adds a fifth operator $E$ (Entropy / Generative
Time Circuit), whose action variable accumulates the irreversible cost:
\[
  C \;\longrightarrow\; K \;\longrightarrow\; F
  \;\longrightarrow\; U \;\longrightarrow\; E \;\longrightarrow\; C'
  \;\longrightarrow\; \cdots
\]

The central mathematical claim is that this operator sequence is not
an analogy or a metaphor: it is a precise structure that admits
constructive definitions, analytical invariants, a variational
principle, and a Hamiltonian formulation.
The framework is placed within established mathematical traditions:
comparison geometry, singularity theory, geometric flows, variational
mechanics, and impulsive Hamiltonian systems.

The paper is self-contained.
Section~\ref{sec:assumptions} states the six minimal assumptions.
Section~\ref{sec:operators} gives constructive operator definitions.
Section~\ref{sec:falsifiability} states four falsifiability conditions.
Section~\ref{sec:theorems} proves five structural theorems.
Section~\ref{sec:examples} gives canonical examples.
Section~\ref{sec:invariants} lists seven analytical invariants.
Section~\ref{sec:normalforms} gives the four normal forms.
Section~\ref{sec:singularities} classifies admissible singularities.
Section~\ref{sec:kappa} develops the metric geometry of $\kappa^*$.
Section~\ref{sec:variational} derives the variational principle.
Section~\ref{sec:hamiltonian} develops the Hamiltonian structure.
Section~\ref{sec:gcm} connects to the dm$^3$ contact-geometric framework.
Section~\ref{sec:entropy} introduces the fifth operator $E$.
Section~\ref{sec:perelman} establishes the Perelman correspondence.
Section~\ref{sec:threshold} states the dimensional threshold conjecture.
Section~\ref{sec:status} summarises formal status.
Section~\ref{sec:corpus} lists the complete series corpus.

% ==================================================================
\section{Minimal Mathematical Assumptions}
\label{sec:assumptions}
% ==================================================================

The operator sequence is defined under the weakest assumptions
necessary for the constructions of Section~\ref{sec:operators} to
be well-posed.

\begin{assumption}[Manifold Structure]
\label{ass:manifold}
The state space $X$ is a smooth, finite-dimensional Riemannian
manifold, locally compact and second-countable.
\end{assumption}

\begin{assumption}[Trajectory Regularity]
\label{ass:trajectory}
A system trajectory $\gamma : [0,T] \to X$ is piecewise $C^2$,
locally non-degenerate ($\|\dot\gamma(t)\| \neq 0$ a.e.), and has
bounded curvature on compact intervals prior to folding events.
\end{assumption}

\begin{assumption}[Compression Feasibility]
\label{ass:compression}
There exists a lower-dimensional submanifold $X_C \subset X$ and a
Lipschitz projection $C : X \to X_C$ satisfying the non-collapse
condition
\[
  d\bigl(C(x_1), C(x_2)\bigr) \;\geq\; \delta\, d(x_1, x_2)
\]
for some $\delta > 0$ and all $x_1, x_2$ in a compact neighbourhood.
This is a bi-Lipschitz embedding condition ensuring distinct
trajectories remain distinguishable after compression.
\end{assumption}

\begin{assumption}[Curvature Threshold]
\label{ass:kappa}
The critical curvature is defined intrinsically by the focal radius:
\[
  \kappa^*(x) = \frac{1}{\mathrm{foc}(x)},
\]
where $\mathrm{foc}(x)$ is the distance from $x$ to the first focal
point along the normal direction.
In the presence of positive sectional curvature $K_\mathrm{sec}(x) > 0$,
this is modified by the Rauch comparison theorem to
\[
  \kappa^*(x) = \min\!\bigl(\|\mathrm{II}_x\|,
    \sqrt{K_\mathrm{sec}(x)}\bigr),
\]
where $\|\mathrm{II}_x\|$ is the norm of the second fundamental form.
For $K_\mathrm{sec}(x) \leq 0$, $\kappa^*(x) = \|\mathrm{II}_x\|$.
\end{assumption}

\begin{assumption}[Folding Well-Posedness]
\label{ass:fold}
The folding operator $F : X_C \to X_F$ satisfies: the Jacobian $dF$
loses rank by exactly 1 at fold points; the fold is local; and the
fold produces a finite number of branches.
\end{assumption}

\begin{assumption}[Morse Stability Functional]
\label{ass:morse}
The stability functional $\Phi : X \to \mathbb{R}$ is $C^2$, bounded
below on compact subsets, and Morse: all critical points are
non-degenerate, i.e., $\nabla^2\Phi(x^*) \succ 0$ at every local
minimum $x^*$.
\end{assumption}

% ==================================================================
\section{Operator Definitions}
\label{sec:operators}
% ==================================================================

Let $X$ be as in Assumption~\ref{ass:manifold}.
A system trajectory is a piecewise $C^2$ map
$\gamma : [0,T] \to X$.
The goal is to describe local transitions in $\gamma$ via the sequence
$C \to K \to F \to U$.

\begin{definition}[Compression Operator $C$]
\label{def:C}
A \emph{compression operator} is a map $C : X \to X_C$ with
$\dim(X_C) < \dim(X)$ satisfying Assumption~\ref{ass:compression}.
$C$ reduces degrees of freedom while retaining essential local
structure.
\end{definition}

\begin{definition}[Curvature Operator $K$]
\label{def:K}
Given a compressed trajectory $\gamma_C : [0,T] \to X_C$,
the \emph{curvature operator} $K : X_C \to X_C$ modifies the tangent
field by
\[
  \frac{d}{ds}\bigl[K(\gamma_C)(s)\bigr]
  = \dot\gamma_C(s) + \alpha(s)\,\mathbf{n}(s),
\]
where $\mathbf{n}(s)$ is the unit normal and
$\alpha(s) = \lambda\bigl(\kappa^*(\gamma_C(s)) - \kappa(s)\bigr)_+$,
$\lambda > 0$.
$K$ is a curvature-inducing flow that drives curvature monotonically
toward $\kappa^*$ but never beyond it; the sequence is a hybrid
dynamical system with a switching condition at $\kappa^*$.
\end{definition}

\begin{definition}[Folding Operator $F$]
\label{def:F}
Let $\gamma_K = K(\gamma_C)$.
A fold occurs at $s_0$ when $|\kappa_K(s_0)| = \kappa^*(\gamma_K(s_0))$.
The \emph{folding operator} $F : X_C \to X_F$ acts by
\[
  F(\gamma_K(s)) = \gamma_K(s) - \beta(s)\,\mathbf{n}(s),
  \qquad
  \beta(s) =
  \begin{cases}
    0 & |\kappa_K(s)| < \kappa^*, \\
    \mu\bigl(|\kappa_K(s)| - \kappa^*\bigr) & |\kappa_K(s)| \geq \kappa^*,
  \end{cases}
\]
where $\mu > 0$ controls fold depth.
At a fold point $s_0$, $\mathrm{rank}(dF_{\gamma_K(s_0)}) =
\dim(X_C) - 1$.
\end{definition}

\begin{definition}[Unfolding Operator $U$]
\label{def:U}
The \emph{unfolding operator} $U : X_F \to X$ is
\[
  U(x_F) = \operatorname*{arg\,min}_{y \in \mathcal{N}(x_F)} \Phi(y),
\]
where $\mathcal{N}(x_F)$ is a sufficiently small neighbourhood of
$x_F$ and $\Phi$ satisfies Assumption~\ref{ass:morse}.
The unfolding is realised by the gradient flow
$\dot{y} = -\nabla\Phi(y)$, $y(0) = x_F$,
converging to a non-degenerate local minimum $x^*$.
\end{definition}

\begin{theorem}[Sequential Consistency]
\label{thm:consistency}
If $K$ is applied until curvature reaches $\kappa^*$, then $F$ is
well-defined, produces a finite branch set, and induces a rank-deficient
Jacobian at fold points.
\end{theorem}

\begin{proof}[Proof sketch]
$K$ drives curvature monotonically to $\kappa^*$ via the gain field
$\alpha(s)$.
At $\kappa^*$, the term $\beta(s)$ becomes nonzero, introducing local
non-injectivity.
The normal direction collapses, reducing Jacobian rank by~1.
Locality and the finite critical-point structure of $\kappa^*$ (via the
Morse assumption on $\Phi$) imply finitely many branches.
\end{proof}

% ==================================================================
\section{Falsifiability Conditions}
\label{sec:falsifiability}
% ==================================================================

The framework is falsifiable at four levels.

\textbf{F1 --- Compression.}
If empirical data show expansion under $C$, or collapse of distinct
trajectories, the model fails.

\textbf{F2 --- Curvature.}
If a fold occurs at curvature strictly below $\kappa^*$, or curvature
exceeds $\kappa^*$ without folding, the model is falsified.
($\kappa^*$ is computed independently via the focal radius, not
post-hoc from fold observations.)

\textbf{F3 --- Folding.}
If empirical fold events do not correspond to Jacobian rank loss,
or produce infinitely many branches, the model fails.

\textbf{F4 --- Sequence order.}
If transitions occur in a different order, or stabilisation occurs
without folding, the model is falsified.

% ==================================================================
\section{Structural Theorems}
\label{sec:theorems}
% ==================================================================

\begin{theorem}[Existence and Well-Posedness]
\label{thm:existence}
Under Assumptions~\ref{ass:manifold}--\ref{ass:morse}, the composite
operator $G = U \circ F \circ K \circ C$ is well-defined on any
piecewise $C^2$ trajectory.
\end{theorem}

\begin{theorem}[Local Determination]
\label{thm:local}
The action of $G$ on $\gamma$ is determined entirely by the local
geometry of $\gamma$ in a neighbourhood of the fold point.
\end{theorem}

\begin{theorem}[Non-Commutativity]
\label{thm:noncomm}
The operators $C, K, F, U$ do not commute; the sequence is
order-dependent.
\end{theorem}

\begin{theorem}[Irreducibility]
\label{thm:irred}
No operator in the sequence $C \to K \to F \to U$ can be removed
without altering the qualitative structure of the transition.
\end{theorem}

\begin{theorem}[Finite Branching]
\label{thm:finite}
The branch set
$\mathcal{B} = \{F(\gamma_K(s_i)) : |\kappa_K(s_i)| = \kappa^*(\gamma_K(s_i))\}$
is finite.
\end{theorem}

\begin{remark}
Finiteness follows from the Morse condition on $\Phi$
(Assumption~\ref{ass:morse}), which prevents fourth-order degeneracies
and ensures isolated critical points, together with the transversality
of $\gamma$ to the fold locus (a generic condition assumed explicitly
in applications).
\end{remark}

% ==================================================================
\section{Canonical Examples}
\label{sec:examples}
% ==================================================================

\textbf{Planar curve with curvature-driven flow.}
$\gamma : [0,T] \to \mathbb{R}^2$, with $C$ an orthogonal projection,
$K$ the curvature-inducing flow, $F$ activated at $\kappa^*$, and $U$
gradient descent on a local $\Phi$.
The simplest non-trivial realisation of the sequence.

\textbf{Elastic rod under compression.}
An elastic rod minimising Euler--Bernoulli energy
$E[\gamma] = \int \kappa(s)^2\,ds$.
Axial loading provides $C$; buckling provides $K$; the critical load
is $\kappa^*$; post-buckling configuration provides $U$.
This is the cleanest physical realisation of the curvature threshold.

\textbf{Gradient flow on a double-well potential.}
$\Phi(x) = (x^2 - 1)^2$.
The unstable equilibrium at $x = 0$ is the fold point; gradient flow
selects one of $x = \pm 1$.
Illustrates finite branching and stability selection.

\textbf{Saddle-node bifurcation.}
$\dot{x} = \mu - x^2$, $\dot{y} = -y$.
Projection onto the slow manifold provides $C$; approach to the fold
provides $K$; loss of the slow manifold at $\mu = 0$ provides $F$;
flow to the stable branch provides $U$.

\textbf{dm$^3$ toy model.}
The canonical contact-geometric realisation on
$M = S^1 \times \mathbb{R}^+ \times \mathbb{R}$:
$\dot{r} = r(1-r^2) + 2(r-1)e^{-z}$,
$\dot\theta = 1$,
$\dot{z} = r^2 - 2(r-1)^2 e^{-z}$.
The limit cycle $\Gamma$ at $r = 1$ is the post-fold stabilised state;
the entropy operator $E$ is $\dot{z}|_\Gamma = 1 > 0$.
See Section~\ref{sec:entropy} and Figure~\ref{fig:phase}.

% ==================================================================
\section{Analytical Invariants}
\label{sec:invariants}
% ==================================================================

\begin{definition}[Analytical Invariants]
The following quantities are preserved by the operator sequence $G$:

\begin{enumerate}[label=\textbf{I\arabic*.},leftmargin=3em]
  \item \textbf{Ambient dimension.}
    $\dim(X)$ is preserved by $G$.
    No claim is made about the dimension of the image of a specific
    trajectory.

  \item \textbf{Codimension of the fold.}
    $\mathrm{codim}(F(\gamma)) = 1$, following from the rank-loss
    condition.

  \item \textbf{Critical curvature threshold.}
    $\kappa^*(x)$ is a geometric property of the manifold, invariant
    under the operator sequence.

  \item \textbf{Curvature sign.}
    $\mathrm{sgn}(\kappa(s))$ is preserved under $K$ and $F$; only
    magnitude changes.

  \item \textbf{Injectivity before threshold.}
    For all $s$ with $|\kappa(s)| < \kappa^*(s)$, the trajectory
    remains injective under $C$, $K$, $F$.

  \item \textbf{Rank deficiency at fold.}
    $\mathrm{rank}(dF) = \dim(X_C) - 1$ at every fold point.

  \item \textbf{Energy monotonicity under $U$.}
    $\Phi(U(x)) \leq \Phi(x)$, with strict inequality unless $x$ is
    already a local minimum.
\end{enumerate}
\end{definition}

\begin{remark}
An energy-gap invariant $\Delta\Phi = \Phi(x_2) - \Phi(x_1)$ is
preserved under $C, K, F$ only if $\Phi$ is constant along the fibers
of each operator.
This is not guaranteed in general and is not claimed as a universal
invariant.
\end{remark}

% ==================================================================
\section{Normal Forms}
\label{sec:normalforms}
% ==================================================================

Two transitions $G_1, G_2$ are \emph{equivalent} if there exists a
diffeomorphism $\psi : X \to X$ with $G_2 = \psi \circ G_1 \circ
\psi^{-1}$.

\begin{definition}[Normal Forms]
\label{def:normalforms}
\begin{enumerate}[label=\textbf{NF\arabic*.},leftmargin=3em]
  \item \textbf{Compression.}
    $C_\mathrm{NF} = \pi_k : \mathbb{R}^n \to \mathbb{R}^k$,
    the standard coordinate projection.

  \item \textbf{Curvature.}
    $K_\mathrm{NF}(s) = (s, \alpha s^2)$, $\alpha > 0$.
    This produces curvature $\kappa(s) = 2\alpha(1 + 4\alpha^2 s^2)^{-3/2}$,
    approximating $2\alpha$ near $s = 0$.

  \item \textbf{Folding --- Whitney fold.}
    $F_\mathrm{NF}(u,v) = (u, v^2)$.
    This is the unique (up to diffeomorphism) normal form for a
    codimension-1 fold singularity \cite{Arnold1986}.

  \item \textbf{Unfolding.}
    $U_\mathrm{NF}(x) = 0$, the canonical stabilisation map arising
    from $\Phi_\mathrm{NF}(x) = x^2$.
\end{enumerate}
\end{definition}

% ==================================================================
\section{Singularity Classification}
\label{sec:singularities}
% ==================================================================

\subsection{Admissible Singularities}

Given the constraints of rank-1 loss, finite branching, and the Morse
condition on $\Phi$, the admissible singularities are:
\begin{enumerate}[leftmargin=2em]
  \item Fold $A_1$: $(u,v) \mapsto (u, v^2)$
  \item Cusp $A_2$: $(u,v) \mapsto (u, v^3 + uv)$
  \item Swallowtail $A_3$: $(u,v) \mapsto (u, v^4 + uv^2 + \beta v)$
\end{enumerate}
Higher $A_k$ ($k \geq 4$) require a $k$-parameter unfolding.
The Morse condition on $\Phi$ restricts the unfolding to at most three
parameters, excluding $A_4$ and above.

\subsection{Classification Conditions}

Let $\Delta(s) = \kappa(s) - \kappa^*(s)$.

\begin{center}
\small
\begin{tabular}{llll}
\toprule
Type & Conditions & Normal form \\
\midrule
$A_1$ (fold)        & $\Delta(s_0)=0$, $\Delta'(s_0)\neq 0$
                    & $(u, v^2)$ \\
$A_2$ (cusp)        & $\Delta=\Delta'=0$, $\Delta''\neq 0$
                    & $(u, v^3+uv)$ \\
$A_3$ (swallowtail) & $\Delta=\Delta'=\Delta''=0$, $\Delta'''\neq 0$
                    & $(u, v^4+uv^2+\beta v)$ \\
\bottomrule
\end{tabular}
\end{center}

\begin{theorem}[Classification of Generative Transitions]
\label{thm:classification}
Every admissible generative transition $G = U \circ F \circ K \circ C$
is $\mathcal{A}$-equivalent to exactly one of $A_1$, $A_2$, $A_3$.
\end{theorem}

\begin{proof}[Proof sketch]
Rank loss is exactly 1, restricting to the $A_k$ series.
The Morse condition on $\Phi$ limits the unfolding to at most three
parameters.
Transversality of $\gamma$ to the fold locus (generic assumption)
ensures isolated fold points.
Thus $k \leq 3$.
\end{proof}

\subsection{Bifurcation Stratification}

Parameter space $\Theta \subset \mathbb{R}^p$, $p \leq 3$,
decomposes into:
\begin{itemize}[leftmargin=2em]
  \item $\Theta_1$: generic fold region (codimension 0 in $\Theta$),
  \item $\Theta_2$: cusp stratum (codimension 1),
  \item $\Theta_3$: swallowtail stratum (codimension 2; a single point
    when $p = 3$).
\end{itemize}

% ==================================================================
\section{Metric Geometry of $\kappa^*$}
\label{sec:kappa}
% ==================================================================

\subsection{Intrinsic Definition}

$\kappa^*(x) = 1/\mathrm{foc}(x)$, where $\mathrm{foc}(x)$ is the
focal radius.
For a curve embedded in $(X, g)$, this equals $\|\mathrm{II}_x\|$
when $K_\mathrm{sec} \leq 0$.

\subsection{Sectional Curvature Correction}

By the Rauch comparison theorem, for variable positive sectional
curvature:
\[
  \kappa^*(x) = \min\!\bigl(\|\mathrm{II}_x\|,
    \sqrt{K_\mathrm{sec}(x)}\bigr).
\]
Positive ambient curvature lowers the threshold for folding.

\subsection{Computability}

Given $\gamma$ in $(X, g)$:
(1) compute $\|\mathrm{II}_x\|$ from the second fundamental form;
(2) compute $K_\mathrm{sec}(x)$ from the curvature tensor;
(3) apply the formula above.
The threshold is stable:
$\delta\kappa^* = O(\delta g) + O(\delta\mathrm{II}) +
O(\delta K_\mathrm{sec})$.

% ==================================================================
\section{Variational Principles}
\label{sec:variational}
% ==================================================================

Define the action functional
\[
  S[\gamma] = \int_0^T \!\Bigl[
    \tfrac{1}{2}\|P^\perp \dot\gamma\|^2
    + \tfrac{\lambda}{2}(\kappa^* - \kappa)_+^2
    + \mu\,\delta(|\kappa| - \kappa^*)
    + \Phi(\gamma)
  \Bigr]\,ds.
\]
Each term corresponds to one operator:
\begin{itemize}[leftmargin=2em]
  \item $L_C = \tfrac{1}{2}\|P^\perp\dot\gamma\|^2$:
    minimise orthogonal kinetic energy (compression).
  \item $L_K = \tfrac{\lambda}{2}(\kappa^* - \kappa)_+^2$:
    minimise curvature deficit.
  \item $L_F = \mu\,\delta(|\kappa| - \kappa^*)$:
    singular activation at threshold.
  \item $L_U = \Phi(\gamma)$: minimise stability potential.
\end{itemize}

The delta term places the framework in the class of free-discontinuity
variational problems (cf.\ Ambrosio--Tortorelli, Mumford--Shah,
Francfort--Marigo).
It produces the jump condition
\[
  \Bigl[\frac{\partial L}{\partial\dot\gamma}\Bigr]_{s_0^-}^{s_0^+}
  = \mu\,\mathbf{n}(s_0),
\]
which is the variational counterpart of the fold map.
The full generative transition is therefore a piecewise-smooth
extremal of $S$.

% ==================================================================
\section{Hamiltonian Discontinuities and Symplectic Geometry}
\label{sec:hamiltonian}
% ==================================================================

\subsection{Phase Space}

The phase space is $T^*X$ with canonical coordinates $(\gamma, p)$
and symplectic form $\omega = d\gamma \wedge dp$.

\subsection{Smooth Flow}

For $s \neq s_0$, the system obeys Hamilton's equations and the flow
is symplectic: $\Phi_t^*\omega = \omega$.

\subsection{Momentum Jump at the Fold}

The delta Lagrangian produces the impulse
\[
  p(s_0^+) - p(s_0^-) = \mu\,\mathbf{n}(s_0).
\]
Configuration $\gamma$ is continuous; momentum $p$ has a jump.

\subsection{The Fold as a Symplectic Map}

Define the fold map in phase space:
$\mathcal{F} : (\gamma, p) \mapsto (\gamma, p + \mu\mathbf{n})$.

\begin{theorem}[Symplectic Preservation]
\label{thm:symplectic}
$\mathcal{F}^*\omega = \omega$.
\end{theorem}

\begin{proof}
$d\gamma \wedge d(p + \mu\mathbf{n}) = d\gamma \wedge dp +
\mu\,d\gamma \wedge d\mathbf{n}$.
Since $\mathbf{n}$ depends only on $\gamma$, the term
$d\gamma \wedge d\mathbf{n} = 0$.
Hence $\mathcal{F}^*\omega = \omega$.
\end{proof}

\subsection{Canonical Transformation}

The fold is generated by the distributional function
$S(\gamma) = \mu\,\Theta(|\kappa(\gamma)| - \kappa^*)$,
so that $p^+ = p^- + \partial S/\partial\gamma$.

\subsection{Full Transition}

Let $\Phi_t$ be the pre-fold Hamiltonian flow, $\mathcal{F}$ the fold
map, $\Psi_t$ the post-fold flow.
The consistency condition (basin structure) ensures the output of
$\mathcal{F}$ lies in the domain of $\Psi_t$.
Then $H = \Psi_t \circ \mathcal{F} \circ \Phi_t$ is a
piecewise-smooth symplectic map: $H^*\omega = \omega$.

\subsection{Singularity Type and Momentum Structure}

\begin{itemize}[leftmargin=2em]
  \item $A_1$ fold: single momentum jump.
  \item $A_2$ cusp: momentum jump with first-order tangency constraint.
  \item $A_3$ swallowtail: momentum jump with second-order tangency.
\end{itemize}

% ==================================================================
\section{Connection to the dm$^3$ Framework}
\label{sec:gcm}
% ==================================================================

This section makes precise the relationship between the
singularity-theoretic framework of this volume and the
contact-geometric dm$^3$ framework of the companion papers
\cite{Grossi2026_SIAM, Grossi2026_ChA}.

\subsection{From Trajectories to Limit Cycles}

In this volume the central object is a trajectory
$\gamma : [0,T] \to (X,g)$ undergoing a localised generative
transition.
In the dm$^3$ framework, the central object is a hyperbolic limit cycle
$\Gamma \subset M$ embedded in a dissipative dynamical system.
The link: every admissible generative transition induces a locally
attracting invariant set; conversely, every dm$^3$ limit cycle admits
a neighbourhood whose formation is governed by a fold-type transition.

\subsection{Curvature Threshold vs.\ Embodiment Threshold}

In this volume, folding is triggered at the geometric threshold
$\kappa^*(x) = 1/\mathrm{foc}(x)$.
In the dm$^3$ framework, stochastic stability is governed by the
embodiment threshold $\tau = \sqrt{c/\kappa}$ defined by the stochastic
Lyapunov condition $\mathcal{L}V \leq -cV + \kappa\|\sigma\|^2$.
As $\kappa \uparrow \kappa^*$, the folding operator $F$ introduces
rank-1 loss and the unfolding $U$ selects a stable branch; the
resulting orbit $\Gamma$ has Floquet exponent $\mu_{\max} < 0$.
By the stochastic Lyapunov condition, $\tau = \sqrt{c/\kappa}$ is
finite precisely when $\mu_{\max} < 0$.
Thus $\kappa^*$ is the geometric precursor of $\tau$.

\subsection{Operator Grammar Correspondence}

\begin{center}
\small
\begin{tabular}{ll}
\toprule
This volume & dm$^3$ framework \\
\midrule
Compression $C$ & Basin contraction \\
Curvature flow $K$ & Lyapunov descent $\dot{V} \leq -cV$ \\
Fold $F$ & Whitney $A_1$ at $q^* = 1$; contact bifurcation \\
Unfolding $U$ & Gradient flow to $\Gamma$ \\
Entropy $E$ (second edition) & Entropy operator $\dot{z} \geq 0$ \\
\bottomrule
\end{tabular}
\end{center}

\subsection{Singularity Classification vs.\ Bifurcation Diagram}

This volume proves that admissible generative transitions are precisely
$A_1$, $A_2$, $A_3$.
The dm$^3$ framework identifies exactly four bifurcations: contact
Hopf, saddle-node of limit cycles, Neimark--Sacker, and slow-fast
crossover.
The $A_1$--$A_3$ singularities classify the local geometry of these
bifurcations under the projection $M = S \times \mathbb{R} \to S$:
$A_1$ corresponds to the contact Hopf and saddle-node;
$A_2$ to the Neimark--Sacker;
$A_3$ to the slow-fast crossover.

\subsection{Summary of the Three Papers}

\begin{itemize}[leftmargin=2em]
  \item \textbf{This volume:} explains why folds must occur and
    classifies their local geometry.
  \item \textbf{dm$^3$ toy model \cite{Grossi2026_SIAM}:} proves that
    the resulting dynamics are computable, stable, and structurally
    closed.
  \item \textbf{Chapter A \cite{Grossi2026_ChA}:} applies the
    framework to autophagy and the triple-alpha process; 18 theorems
    machine-checked in Lean~4.
\end{itemize}

% ==================================================================
\section{The Fifth Operator: Entropy and the Generative Time Circuit}
\label{sec:entropy}
% ==================================================================

\subsection{The Action Variable as Entropy}

The contact manifold $M = S \times \mathbb{R}$ carries a coordinate
$z \in \mathbb{R}$ satisfying $\dot{z} = f(x) \geq 0$ near the limit
cycle $\Gamma$.
In the thermodynamic interpretation of the contact structure, the form
$\alpha = dz - \lambda$ encodes both the first and second laws: $dz$
is the entropy production, $\lambda$ is the reversible work.
The condition $\dot{z} \geq 0$ is the second law on $M$.

\begin{definition}[Generative Time Circuit Operator $E$]
\label{def:E}
The \emph{Generative Time Circuit} operator
$E : \Gamma \to \Gamma'$ maps the limit cycle of one generative
cycle to the compression parameter of the next:
\[
  E : z(T^*) \;\mapsto\; \kappa' = \phi(z(T^*)),
\]
where $\phi' < 0$ (accumulated entropy increases compression in the
next cycle) and $\phi(0) = \kappa_0$.
\end{definition}

The chain closes as a spiral:
$C \to K \to F \to U \to E \to C' \to \cdots$

\begin{figure}[ht]
\centering
\includegraphics[width=\linewidth]{fig_A2_phase_portrait}
\caption{dm$^3$ phase portrait in contact coordinates $(\rho, \theta)$.
\textbf{Left:} autophagy manifold $X_\mathrm{auto}$.
\textbf{Right:} stellar interior manifold $X_\mathrm{star}$.
Gold circle: limit cycle $\Gamma$ at $r = 1$.
Dashed circles: Gronwall basin $\varepsilon_0 = 1/3$ and numerical
inner boundary $r^* \approx 0.80$.
Blue: converging orbits; red: escaping orbits.
The two portraits are identical up to axis relabelling, demonstrating
the contact morphism.
Lean: \texttt{gronwall\_radius}, \texttt{basin\_asymmetry}.}
\label{fig:phase}
\end{figure}

\begin{theorem}[T1 --- Entropy Monotonicity; \emph{proof obligation}]
\label{thm:T1}
Along trajectories of the dm$^3$ toy model in the Gronwall basin,
$z(t)$ is monotonically non-decreasing for all $t > 0$, and
$\dot{z}|_\Gamma = 1 > 0$.
\end{theorem}

\begin{remark}
The scalar claim $\dot{z}|_\Gamma = 1 > 0$ is immediate from the toy
ODE at $r = 1$, $z \to \infty$.
Global monotonicity in the basin is an open Lean~4 obligation
(AXLE Issue~\#15).
\end{remark}

% ==================================================================
\section{The Perelman Structural Correspondence}
\label{sec:perelman}
% ==================================================================

Perelman's proof of the Poincar\'e conjecture
\cite{Perelman2002,Perelman2003a,Perelman2003b} proceeds through
Ricci flow with surgery, the $\mathcal{F}$-functional, and the
$\mathcal{W}$-entropy.
The dm$^3$ framework identifies a term-by-term structural
correspondence between the five-operator chain and Perelman's proof.
This correspondence does not re-prove the Poincar\'e conjecture;
Perelman's proof stands independently.

\begin{figure}[ht]
\centering
\includegraphics[width=\linewidth]{fig_1A1_perelman_correspondence}
\caption{The dm$^3$ five-operator chain mapped onto Perelman's Ricci
flow with surgery. Each operator is colour-coded; gold arrows indicate
the structural mapping. Lean~4 status: operators $C,K,F,U$ are
verified (Theorems~A--D); operator $E$ carries an open obligation
(AXLE Issue~\#15). The correspondence is structural, not causal.}
\label{fig:perelman}
\end{figure}

\begin{table}[ht]
\centering
\small
\caption{Term-by-term structural correspondence between the dm$^3$
five-operator chain and Perelman's Ricci flow with surgery.}
\label{tab:perelman}
\begin{tabularx}{\textwidth}{lXX}
\toprule
Op. & dm$^3$ Role & Perelman / Ricci Flow Role \\
\midrule
$C$ & Compression; drives $\kappa \to \kappa^*$ &
  Initial metric selection \\
$K$ & Curvature intensification &
  $\partial_t g_{ij} = -2\,\mathrm{Ric}_{ij}$; diffusion to constant
  curvature \\
$F$ & Whitney $A_1$ at $\kappa^*$; Jacobian rank loss &
  Singularity formation; surgery excises necks \\
$U$ & Stabilisation on limit cycle $\Gamma$ &
  Post-surgery convergence to round metric $S^3$ \\
$E$ & $\dot{z} \geq 0$; seeds next cycle &
  $\mathcal{W}$-entropy; monotone; governs singularity control \\
\bottomrule
\end{tabularx}
\end{table}

\begin{conjecture}[Perelman Structural Correspondence]
\label{conj:perelman}
There exists a functor
$\mathcal{P} : \mathbf{dm^3} \to \mathbf{RicciFlow}$
that maps the compression operator $C$ to initial metric selection;
$K$ to the Ricci tensor acting as diffusion;
$F$ to surgical excision at $\kappa^*$-necks;
$U$ to post-surgery convergence to the round metric;
and $E$ to the $\mathcal{W}$-entropy functional.
$\mathcal{P}$ preserves the chain ordering and maps the dm$^3$ limit
cycle $\Gamma$ to $S^3$ as the terminal attractor.
\end{conjecture}

\begin{remark}
Conjecture~\ref{conj:perelman} is argued but not proved.
The structural parallel is term-by-term (Table~\ref{tab:perelman}).
The functor construction requires: (a) contact morphisms as morphisms
in $\mathbf{dm^3}$ (already defined); (b) surgery-compatible
diffeomorphisms as morphisms in $\mathbf{RicciFlow}$ (standard);
(c) construction of $\mathcal{P}$ and verification of functor laws
(open obligation, DM3-lab \cite{DM3lab2026}).
\end{remark}

% ==================================================================
\section{The Dimensional Threshold: $N = 3$ and $c = 3$}
\label{sec:threshold}
% ==================================================================

\subsection{Contact Geometry and Dimension 3}

In dimension 1, contact structure is trivial.
In dimension 2, Liouville's theorem rules out limit cycle attractors
in area-preserving flows.
In dimension 3, the contact form $\alpha$ first admits a Reeb vector
field with a non-trivial flow and a limit cycle attractor
\cite{Geiges2008}.
Dimension 3 is the \emph{minimum dimension for non-trivial contact
geometry}.
This is why the dm$^3$ framework lives on a 3-dimensional contact
manifold: not by convention but by structural necessity.
It is also why the Poincar\'e conjecture required Perelman's
Ricci flow: dimensions $\geq 5$ (Smale, 1961 \cite{Smale1961})
and $4$ (Freedman, 1982 \cite{Freedman1982}) yield to other methods;
dimension 3, the minimum for non-trivial contact geometry, is the
hardest case.

\subsection{The Collatz Connection}

\begin{conjecture}[Dimensional Threshold]
\label{conj:threshold}
The constant $c = 3$ in the Collatz map $n \mapsto 3n+1$ and the
dimension $N = 3$ in the Poincar\'e conjecture are both instances of
the same abstract threshold: the minimum value at which a generative
system first admits non-trivial contact-geometric structure.
$c = 3$ is the minimum value for which the discrete analogue of the
dm$^3$ operator chain generates a non-collapsing, non-periodic orbit
structure analogous to a hyperbolic limit cycle.
\end{conjecture}

\begin{remark}
Conjecture~\ref{conj:threshold} is motivated but not proved.
Formal proof requires a definition of contact-geometric structure for
discrete dynamical systems on $\mathbb{Z}$ and a proof that the
dm$^3$ axioms have no discrete analogue for $c < 3$ (open obligation,
DM3-lab).
\end{remark}

% ==================================================================
\section{Formal Status and Open Obligations}
\label{sec:status}
% ==================================================================

\subsection{What Is Machine-Checked (Lean~4, Zero \texttt{sorry})}

\begin{enumerate}[label=\textbf{P\arabic*.},leftmargin=3em]
  \item Whitney $A_1$ conditions on $V(q) = q^3-3q$ at $q=1$:
    $V'(1)=0$, $V''(1)=6\neq 0$, $V(1)=-2$,
    $V(q)+2=(q-1)^2(q+2)$.
  \item Contact non-degeneracy: $c(\rho)=-2\rho < 0$ for $\rho > 0$.
  \item Gronwall radius: $\varepsilon_0 = 1/3$.
  \item Basin asymmetry: $1/3 < 4/5 \approx r^*$.
  \item Lyapunov exponents: $-V''(1)/2=-3$; $\mu_{\max}=-2 < 0$.
  \item Stability functional: $\sigma(\rho)=\rho^2 > 0$,
    $\sigma'(\rho) > 0$ for $\rho > 0$.
  \item Theorems A--D (existence, closure, normal form, stability).
\end{enumerate}

\begin{figure}[ht]
\centering
\includegraphics[width=\linewidth]{fig_A4_coherence_bridge}
\caption{Coherence Bridge: dm$^3$ scalar invariants $(\mu_{\max}, \beta)$
across all domains.
\textbf{Left:} $\mu_{\max} < 0$ for all domains
(Lean: \texttt{mu\_dm3\_neg}).
\textbf{Right:} fold sharpness $\beta$.
Gold: autophagy (Chapter~A). Teal: triple-alpha (Chapter~A).
Grey: prior domains.}
\label{fig:bridge}
\end{figure}

\subsection{Open Obligations}

\begin{enumerate}[label=\textbf{O\arabic*.},leftmargin=3em]
  \item \textbf{AXLE Issue \#12} (\texttt{kappa\_lipschitz}).
    Lipschitz continuity of $K$ on configuration manifolds.
  \item \textbf{AXLE Issue \#14} (AutophagyDm3).
    Full contact non-degeneracy on $X_\mathrm{auto}$;
    Whitney fold from mTORC1 kinase data;
    existence of $\Gamma_\mathrm{auto}$.
  \item \textbf{AXLE Issue \#15} (Theorem~T1).
    Global monotonicity of $z(t)$ in the Gronwall basin.
  \item \textbf{Sorry~1} (Discrete dm$^3$).
    Extension to discrete dynamical systems on $\mathbb{Z}$.
  \item \textbf{Conjecture~\ref{conj:perelman}} (Perelman functor).
    Construction of $\mathcal{P}$ and functor law verification.
\end{enumerate}

\subsection{What Is Argued (Not Proved)}

The following are complete mathematical arguments without Lean~4
verification, deposited in DM3-lab \cite{DM3lab2026}:
Perelman structural correspondence
(Table~\ref{tab:perelman}, Conjecture~\ref{conj:perelman});
dimensional threshold
(Conjecture~\ref{conj:threshold});
Kakeya needle problem as the $\tau \to 0$ limit on the rotation
contact manifold;
Tribonacci constant as the dm$^3$ structural growth rate for the
three-operator iteration.
These are proof sketches, not proof claims.

% ==================================================================
\section{The Principia Orthogona Corpus}
\label{sec:corpus}
% ==================================================================

\begin{table}[ht]
\centering
\small
\caption{Principia Orthogona corpus, April--May 2026.
All deposits open access CC BY-NC-ND 4.0.
\emph{Submitted} = under journal review;
\emph{Deposited} = open-access preprint.}
\label{tab:corpus}
\begin{tabularx}{\textwidth}{llX}
\toprule
Zenodo DOI & Date & Title / Status \\
\midrule
19117400 & Mar 17 &
  \textit{Principia Orthogona Vol~I.} \emph{Deposited.} \\
19122168 & Mar 19 &
  \textit{GCM --- Geometric Framework for Dissipative Systems.}
  \emph{Submitted: J.\ Geom.\ Mech.} \\
19162013 & Mar 22 &
  \textit{The G6 Crystal.} \emph{Deposited.} \\
19208015 & Mar 24 &
  \textit{Biological Transitions --- Multi-Agent Realisations.}
  \emph{Deposited.} \\
19379385 & Apr 2 &
  \textit{dm$^3$ Operator --- Toy Model \& Global Analysis.}
  \emph{Submitted: SIAM J.\ Appl.\ Dyn.\ Syst.} \\
19379473 & Apr 2 &
  \textit{Principia Orthogona Vol~II --- Contact Realization.}
  \emph{Deposited.} \\
19431918 & Apr 5 &
  \textit{The Number 33 --- Stability Threshold.}
  \emph{Submitted: Mathematical Intelligencer.} \\
19533363 & Apr 2026 &
  \textit{GTCT Paper (Ring 5, Version 2).} \emph{Deposited.} \\
20221723 & May 2026 &
  \textit{Chapter A: Autophagy and Triple-Alpha.}
  \emph{Deposited; 18 Lean theorems, 0 sorry.} \\
\bottomrule
\end{tabularx}
\end{table}

% ==================================================================
\section*{Acknowledgements}
% ==================================================================

This research was conducted independently through G6~LLC,
Newark, New Jersey.
No institutional funding was received.
The author thanks the Lean~4/Mathlib4 community for maintaining
a publicly available formal verification library, and the Zenodo
infrastructure team for open-access deposit services.

% ==================================================================
\section*{Data and Software Availability}
% ==================================================================

Formal verification: \url{https://github.com/TOTOGT/AXLE}.
Numerical simulation and figure generators:
\url{https://github.com/TOTOGT/DM3-lab}.
All deposits open access:
\href{https://doi.org/10.5281/zenodo.19117400}%
{doi:10.5281/zenodo.19117400} (series root).

% ==================================================================
\begin{thebibliography}{99}
\sloppy

\bibitem{Perelman2002}
G.~Perelman,
``The entropy formula for the Ricci flow and its geometric applications,''
\textit{arXiv:math/0211159} (2002).
\url{https://arxiv.org/abs/math/0211159}

\bibitem{Perelman2003a}
G.~Perelman,
``Ricci flow with surgery on three-manifolds,''
\textit{arXiv:math/0303109} (2003).
\url{https://arxiv.org/abs/math/0303109}

\bibitem{Perelman2003b}
G.~Perelman,
``Finite extinction time for the solutions to the Ricci flow on
certain three-manifolds,''
\textit{arXiv:math/0307245} (2003).
\url{https://arxiv.org/abs/math/0307245}

\bibitem{Smale1961}
S.~Smale,
``Generalized Poincar\'e's conjecture in dimensions greater than four,''
\textit{Ann.\ Math.} \textbf{74}(2), 391--406 (1961).
\href{https://doi.org/10.2307/1970239}{doi:10.2307/1970239}

\bibitem{Freedman1982}
M.~H.~Freedman,
``The topology of four-dimensional manifolds,''
\textit{J.\ Differential Geom.} \textbf{17}(3), 357--453 (1982).
\href{https://doi.org/10.4310/jdg/1214437136}{doi:10.4310/jdg/1214437136}

\bibitem{Geiges2008}
H.~Geiges,
\textit{An Introduction to Contact Topology}.
Cambridge University Press, 2008.

\bibitem{Arnold1986}
V.~I.~Arnold,
\textit{Catastrophe Theory}, 3rd~ed. Springer, 1986.

\bibitem{Thom1975}
R.~Thom,
\textit{Structural Stability and Morphogenesis}.
W.~A.~Benjamin, 1975.

\bibitem{AXLE2026}
P.~Nogueira~Grossi,
\textit{AXLE: Lean~4 Formal Verification Engine for the dm$^3$ Programme}.
\url{https://github.com/TOTOGT/AXLE} (2026).

\bibitem{DM3lab2026}
P.~Nogueira~Grossi,
\textit{DM3-lab: Numerical and Structural Verification}.
\url{https://github.com/TOTOGT/DM3-lab} (2026).

\bibitem{Grossi2026_ChA}
P.~Nogueira~Grossi,
``Self-Regulation: Autophagy and the Triple-Alpha Process as
dm$^3$ Generative Transitions,''
Chapter~A of \textit{Principia Orthogona, Book~3}.
Zenodo:
\href{https://doi.org/10.5281/zenodo.20221723}{10.5281/zenodo.20221723}
(2026).

\bibitem{Grossi2026_SIAM}
P.~Nogueira~Grossi,
``dm$^3$ Operator --- Explicit Toy Model and Global Dynamical Analysis,''
Zenodo:
\href{https://doi.org/10.5281/zenodo.19379385}{10.5281/zenodo.19379385}
(2026). \textit{Submitted: SIAM J.\ Appl.\ Dyn.\ Syst.}

\end{thebibliography}

\end{document}
